\section{Ladders: the structure of one walk}\label{sec:ladders}

This section is about a single walk~$x$, pauses allowed, and a horizon~$T$. It equips $x$ with two frame systems, one for its
upper hops (\emph{class~U}) and one for its lower hops (\emph{class~L}, the same construction for $-x$), and proves Theorem~A
(\cref{thm:A}): for a uniformly random offset of the grid of width~$h_0$, the expected number of lost hops of each class is at
most $6/h_0$ times the number of hops of the class. The other walks and the label enter only in
\cref{sec:label,sec:displays,sec:accounting}.

A frame system gives the walk a \emph{band} $\beta(t)$ at every time: the band of the running maximum of its distance from a
reference cell, its \emph{home}. A hop is cheap if the bands at its ends lie in the same group of the grid and two side
conditions hold; bands that differ by~$\Delta$ lie in different groups for at most $\Delta$ of the $h_0$ offsets
(\cref{lem:grid}). So the homes must make
the band differences, summed over all hops, linear in the number of hops. We treat upper hops; lower hops are the upper hops
of~$-x$ (\cref{lem:ladder}(e)).

\subsection{The decomposition by tops}\label{sec:ladders-tops}

\begin{definition}[Tops, ladders and excursions]\label{def:ladder}
\formal{Decomposition.lean}{upperHops, lowerHops, tops, ladder, height}
\formal{FramesU.lean}{base, ladderHop, ladL, ladD, exc, levels}
Let $x$ be a walk and $T$ a horizon.
\begin{enumerate}[label=(\alph*)]
\item The \emph{upper hops} of~$x$ are its hops $(p,u)$ with $x(p+1)>x(p)$, and its \emph{lower hops} those with
$x(p+1)<x(p)$. An upper hop stays above its cell and a lower hop below it (\cref{lem:hops}).
\item The \emph{top} of an upper hop is its owner. The \emph{ladder} of a top~$\tau$ is the set of upper hops with top~$\tau$,
its \emph{height} $H_\tau$ is the number of these hops, its \emph{base} is $y_\tau=x(\tau)-H_\tau$, and its \emph{levels} are the
integers~$w$ with $y_\tau<w<x(\tau)$.
\item For $y_\tau\le z<x(\tau)$, $(L_\tau(z),D_\tau(z))$ denotes the hop of the ladder of~$\tau$ at the cell~$z$; by
\cref{lem:ladder}(a) there is exactly one.
\item The \emph{canonical excursion} of~$\tau$ is $E_\tau=(L_\tau(y_\tau),D_\tau(y_\tau))$, the hop of its ladder at the base.
\end{enumerate}
\end{definition}

\begin{lemma}[Ladders]\label{lem:ladder}
\formal{Decomposition.lean}{image_cell_ladder, ladder_cell_iff, ladder_eq_of_cell, mem_ladder_hopOwner, eq_of_mem_ladder, ladder_lt, ladder_above, ladder_last_visit, ladder_first_visit, ladder_nested, exists_ladder_outer, top_spec, top_ne_start, top_ne_end, card_upperHops_add_card_tops_le, sum_height_eq, upperHops_neg, upperHops_union_lowerHops}
Let $x$ be a walk, $T$ a horizon, and $\tau$ a top of height~$H$ and base~$y$.
\begin{enumerate}[label=(\alph*)]
\item \emph{(One hop per cell.)} The cells of the hops of the ladder of~$\tau$ are exactly $y,y+1,\dots,x(\tau)-1$, and the
ladder has exactly one hop at each of them. Every upper hop lies in exactly one ladder, that of its top.
\item \emph{(Rungs.)} Let $(p,u)$ be the hop of the ladder at the cell~$z$. Then $p<\tau<u$, the walk is above~$z$ on
$(p,u)$, $p$ is the last visit to~$z$ at or before~$\tau$, and $u$ is the first visit to~$z$ at or after~$\tau$. A higher cell
gives a later start and an earlier end, so $E_\tau$ contains every hop of the ladder.
\item \emph{(Tops.)} $1\le\tau$, $\tau+1<T$, $x(\tau-1)<x(\tau)$ and $x(\tau+1)\le x(\tau)$. A top is neither the start nor
the end of an upper hop.
\item \emph{(Counting.)} The number of upper hops plus the number of tops is at most~$T$, and $\sum_\tau H_\tau$, over all
tops, is the number of upper hops.
\item \emph{(Reflection.)} The upper hops of $-x$ are the lower hops of~$x$, and every hop of~$x$ is upper or lower.
\end{enumerate}
\end{lemma}

\begin{outline}
(a) Every hop of the ladder contains~$\tau$ (\lean{ladder_lt}), so two of them at one cell coincide by
\lean{BoundedNesting.hop_unique} (\lean{ladder_eq_of_cell}). If the ladder has a hop at the cell $z_0$ and $z_0\le z<x(\tau)$,
the last visit to~$z$ before~$\tau$ and the first visit after it (\lean{exists_last_visit}, \lean{exists_first_visit}) form an
upper hop on which $\tau$ is again the first maximum (\lean{exists_ladder_cell}); so the cells form an interval ending at
$x(\tau)-1$. Owners are unique (\lean{mem_ladder_hopOwner}, \lean{eq_of_mem_ladder}). (b) A visit to~$z$ in $(p,\tau]$ or
$[\tau,u)$ would lie inside the hop. For $z<z'$ the walk passes~$z'$ between $L_\tau(z)$ and~$\tau$, so $L_\tau(z)<L_\tau(z')$;
symmetrically $D_\tau(z')<D_\tau(z)$. (c) $\tau$ is the first maximum of a hop $(p,u)$, so $x(\tau-1)<x(\tau)$ and
$x(\tau+1)\le x(\tau)$, with $\tau+1\le u<T$; an upper hop starting at~$\tau$ would need $x(\tau+1)>x(\tau)$, and one ending
at~$\tau$ would need $x(\tau-1)>x(\tau)$. (d) Hops with the same start coincide (\lean{BoundedNesting.hop_start_unique}), so the
starts of the upper hops are distinct, and by~(c) they are not tops. The ladders are the fibres of the owner map.
(e) Reflection preserves hops and reverses their first step, which leaves the cell.
\end{outline}

\begin{example}\label{ex:ladders}
On the walk $0,1,2,3,2,3,2,1,0$ with $T=9$, the upper hops are $(0,8)$, $(1,7)$, $(2,4)$ and $(4,6)$, the only lower hop is
$(3,5)$, and the tops are $3$ and~$5$. The ladder of~$3$ is $\{(0,8),(1,7),(2,4)\}$, of height~$3$ and base~$0$; the ladder
of~$5$ is $\{(4,6)\}$, of height~$1$ and base~$2$. The hop $(2,4)$ of the top~$3$ ends where the excursion $E_5=(4,6)$ starts:
an excursion may start at a ladder time of another top. On $0,1,1,0$ the only top is~$1$, and the walk pauses there. These
values are checked by the kernel in \leanfile{Decomposition.lean}.
\end{example}

\paragraph{Why ladders make every hop cheap.} The ladders partition the upper hops, and a top of height~$H$ owns $H$ of them,
one per cell between its base~$y$ and its peak (\cref{lem:ladder}); Theorem~A charges every hop to its ladder. The rung at the
cell $y+d$ starts on the climb, after the walk has risen by~$d$ from the base, and ends on the descent, after the walk has been
at the top. Measured from the base, the band at its start is about $\band(d)$, and the band at its end, a running maximum, is
about $B=\band(H)$. The rung costs at most the sum of the deficits of the two bands below~$B$ (\cref{lem:cost}). The deficit
$B-\band(d)$ exceeds~$j$ only for $d<2^{B-j}\le H/2^j$, so the deficits along the climb sum to $O(H)$ (Lemma~E,
\cref{lem:bandsums}); restarting the descent at geometrically spaced passages adds $O(H)$ more (\cref{lem:perexc}). Between the
rungs the walk is arbitrary and may contain tall excursions; the frames below suspend them, so that they do not disturb the bands
on the ladder.

\subsection{Excursions and structures}\label{sec:ladders-exc}

Let $\tau$ be a top with base~$y$, and $w$ a level of~$\tau$. Between the rungs $L_\tau(w-1)$ and $L_\tau(w)$ of the climb the
walk stays at or above~$w$, and it is at~$w$ at both ends of the \emph{hanging structure} $[L_\tau(w-1)+1,L_\tau(w)]$.
Symmetrically, $[D_\tau(w),D_\tau(w-1)-1]$ is the \emph{descent structure} at level~$w$. A structure is a sequence of upper
excursions above its level, possibly with pauses at the level. The \emph{hanging frame} and the \emph{descent frame} of~$\tau$ at
level~$w$ are
\[
  [L_\tau(w-1)+1,\,L_\tau(w))\qquad\text{and}\qquad[D_\tau(w)+1,\,D_\tau(w-1))
\]
(\lean{hangFrame}, \lean{descFrame}), possibly empty. The descent frame starts one step after its structure, so that the ladder
time $D_\tau(w)$ stays a time of the frame of the ladder. The frames of~$\tau$ are its hanging and descent frames at all its
levels (\lean{framesOf}), and the \emph{suspend frames} of~$x$ are the frames of all its tops (\lean{suspendFrames}).

\begin{lemma}[Excursions and structures]\label{lem:excursions}
\formal{Excursions.lean}{exc_range, upper_laminar, exc_laminar, exc_structure, framesOf_bounds, exists_framesOf_nested, frames_laminar, ladder_time_not_mem_framesOf}
Let $x$ be a walk and $T$ a horizon; let $\tau$ and $\tau'$ be tops, with $E_\tau=(a,b)$ and $E_{\tau'}=(a',b')$. Say that
$E_{\tau'}$ \emph{lies in} $E_\tau$ if $a\le a'$ and $b'\le b$.
\begin{enumerate}[label=(\alph*)]
\item On $[a,b]$ the walk stays in $[y_\tau,x(\tau)]$. Two upper hops are nested or meet in at most one endpoint; in
particular, so are $E_\tau$ and $E_{\tau'}$.
\item \emph{(Structures.)} If $\tau'\ne\tau$ and $E_{\tau'}$ lies in $E_\tau$, then some level~$w$ of~$\tau$ with $w\le
y_{\tau'}$ satisfies either $\tau'<\tau$ and $[a',b']\subseteq[L_\tau(w-1)+1,L_\tau(w)]$, or $\tau<\tau'$ and
$[a',b']\subseteq[D_\tau(w),D_\tau(w-1)-1]$.
\item Every frame $[f_1,f_2)$ of~$\tau$ has $a<f_1$ and $f_2\le b$. In the situation of~(b), $[a'+1,b')$ lies in a frame
of~$\tau$. The root $[0,T)$ and the suspend frames form a laminar family: any two of them are nested or disjoint.
\item The ladder times of~$\tau$, that is, $\tau$ and the times $L_\tau(z)$ and $D_\tau(z)$ for $y_\tau\le z<x(\tau)$, lie in no
frame of~$\tau$ and in no frame of a top whose canonical excursion lies in~$E_\tau$.
\end{enumerate}
\end{lemma}

\begin{outline}
(a) The walk is above~$y_\tau$ inside the outermost ladder hop and at most its maximum $x(\tau)$ (\lean{exc_range}). Upper hops
do not cross (\lean{BoundedNesting.upper_no_cross}), hence are nested or meet in at most an endpoint (\lean{upper_laminar}).
(b) $\tau\notin[a',b']$, since there it would be the first maximum of $E_{\tau'}$ (\lean{top_eq_of_mem}). Say $b'<\tau$. The
time $a'$ is a visit to $y_{\tau'}$ before~$\tau$, so $a'\le L_\tau(y_{\tau'})$, while $L_\tau(y_\tau)=a<a'$ because ladder hops
with a common start have a common top (\lean{top_eq_of_start}). Take the level~$w$ with $L_\tau(w-1)<a'\le L_\tau(w)$
(\lean{exists_switch}); an upper hop that ends by~$\tau$ cannot contain the last visit $L_\tau(w)$ strictly inside
(\lean{not_lt_ladL_lt}), so $b'\le L_\tau(w)$. The case $\tau<a'$ is symmetric (\lean{not_lt_ladD_lt}, \lean{top_eq_of_end}).
(c) The frames of~$\tau$ lie between consecutive rungs inside $(a,b)$ and are pairwise disjoint (\lean{framesOf_disjoint}); by~(b)
the interior of a nested excursion lies in one of them, so frames of nested excursions are nested or disjoint
(\lean{framesOf_laminar_of_nested}), and frames of disjoint excursions are disjoint. (d) A frame of~$\tau$ lies strictly between
consecutive rungs. A frame of another top whose excursion lies in~$E_\tau$ lies strictly inside that excursion, which does not
contain~$\tau$ by~(b), nor any $L_\tau(z)$ or $D_\tau(z)$ strictly inside (\lean{not_lt_ladL_lt}, \lean{not_lt_ladD_lt}).
\end{outline}

The excursions thus form a forest under inclusion, and a nested excursion lies in a single structure of its parent. In
\cref{ex:ladders}, $E_5=(4,6)$ starts at $D_3(2)=4$, a ladder time of its parent, and lies in the descent structure $[4,6]$ of
the top~$3$ at level~$2$, whose descent frame is $[5,7)$: this is why the structures are closed intervals.

\subsection{Frame systems}\label{sec:ladders-framesystems}

\begin{definition}[Frame systems]\label{def:framesystem}
\formal{FrameSystem.lean}{FrameSystem, FrameSystem.Own, FrameSystem.inner, FrameSystem.segStart, FrameSystem.runMax, FrameSystem.home, FrameSystem.beta}
A \emph{frame system}~$\mathcal S$ consists of a horizon~$T$, a map $x\colon\N\to\Z$, a finite set of \emph{frames}, pairs
$F=(f_1,f_2)$ with $f_2\le T$ read as intervals $[f_1,f_2)$, possibly empty, and a finite set~$R$ of \emph{restart points}, such
that the root $[0,T)$ is a frame and any two frames are nested or disjoint.
\begin{enumerate}[label=(\alph*)]
\item A time~$t$ is an \emph{own time} of a frame~$F$ if $F$ contains~$t$ and lies inside every frame that contains~$t$. For
$t<T$, $\inner(t)$ is the frame of which $t$ is an own time (unique by \cref{lem:dynamics}).
\item The \emph{segment start} $s_F(t)$ is the largest element of the set consisting of $f_1$ and of the restart points
$s\in R$ with $f_1<s\le t$ that are own times of~$F$. The pair $(F,s_F(t))$ is the \emph{segment} of~$F$ at~$t$.
\item The \emph{running maximum} $M_F(t)$ is the largest $\abs{x(t')-x(s_F(t))}$ over the own times~$t'$ of~$F$ with
$s_F(t)\le t'\le t$, and $0$ if there is none.
\item The \emph{home} of~$t$ is $\home(t)=x(s_{\inner(t)}(t))$, and its \emph{band} is $\beta(t)=\band(M_{\inner(t)}(t))$.
\end{enumerate}
\end{definition}

\begin{lemma}[Dynamics of a frame system]\label{lem:dynamics}
\formal{FrameDynamics.lean}{FrameSystem.Restart, FrameSystem.segStart_of_restart, FrameSystem.segStart_start, FrameSystem.segStart_succ, FrameSystem.runMax_of_restart, FrameSystem.runMax_start, FrameSystem.runMax_succ_of_not_own, FrameSystem.runMax_succ_of_own, FrameSystem.runMax_frozen, FrameSystem.runMax_le_succ}
\formal{FrameSystem.lean}{FrameSystem.exists_own, FrameSystem.Own.unique, FrameSystem.band_dist_home_le_beta}
Let $\mathcal S$ be a frame system and $F=(f_1,f_2)$ a pair of times. Call~$t$ a \emph{restart} of~$F$ if $t\in R$, $f_1<t$,
and $t$ is an own time of~$F$.
\begin{enumerate}[label=(\alph*)]
\item Every time $t<T$ is an own time of exactly one frame.
\item $s_F(f_1)=f_1$ and $M_F(f_1)=0$. If $t$ is a restart of~$F$, then $s_F(t)=t$ and $M_F(t)=0$.
\item Let $f_1<t$, and let $t$ not be a restart of~$F$. Then $s_F(t)=s_F(t-1)$; moreover $M_F(t)=M_F(t-1)$ if $t$ is not an
own time of~$F$, and $M_F(t)=\max\bigl(M_F(t-1),\abs{x(t)-x(s_F(t))}\bigr)$ if it is.
\item If $f_1\le t_0\le t$ and $F$ has no own time in $(t_0,t]$, then $M_F(t)=M_F(t_0)$. If $f_1<t$ and $t$ is not a
restart of~$F$, then $M_F(t-1)\le M_F(t)$.
\item For $t<T$, $\beta(t)\ge\band\abs{x(t)-\home(t)}$.
\end{enumerate}
\end{lemma}

\begin{outline}
(a) A shortest frame containing~$t$ lies inside every frame containing~$t$, since such frames are nested
(\lean{FrameSystem.nested_of_contains}). (b), (c) The candidate segment starts up to~$t$ are those up to $t-1$, and~$t$ if it is
a restart (\lean{FrameSystem.segStart_eq}); the own times of the current segment up to~$t$ are those up to $t-1$, and~$t$ if it
is an own time. (d) Induction with~(c). (e) $t$ is an own time of $\inner(t)$ in its current segment.
\end{outline}

\paragraph{The hybrid convention.} Only own times enter a running maximum: while a nested frame is innermost, its parent's
running maximum is frozen, and when the nested frame ends, the parent resumes it. Frames are nested, their maxima are not; and a
restart point starts a new segment only in the frame of which it is an own time. The proof uses three consequences. The bands
at the ends of every rung are bounded in terms of the ladder alone, however tall the suspended structures between the rungs are
(\cref{lem:ladderbands}). Consecutive own times of a frame of class~U are at most one level apart, so the maximum of a segment is less than its number of own times, which bounds the weight of the segments as
events of the label (\cref{lem:W1}). And a suspended frame's band is a fixed value that the label can display while the frame is
suspended (\cref{sec:displays}); since every suspend frame ends at band~$0$ (\cref{lem:W2b}), the end of a frame only returns the
walk from band~$0$ to its parent's frozen band, up to the update of the parent's running maximum when it resumes.

\subsection{The frames of class U}\label{sec:ladders-framesU}

\begin{definition}[The frames of class U]\label{def:framesU}
\formal{FramesU.lean}{hangFrame, descFrame, suspendFrames, passages, innerDescents, restartPts}
\formal{Excursions.lean}{framesOf, upperFrames}
Let $x$ be a walk and $T$ a horizon.
\begin{enumerate}[label=(\alph*)]
\item The \emph{suspend frames} are the hanging frames $[L_\tau(w-1)+1,L_\tau(w))$ and the descent frames
$[D_\tau(w)+1,D_\tau(w-1))$ of all tops~$\tau$ at all their levels~$w$.
\item The \emph{passages} of a top~$\tau$ are the times $D_\tau(y_\tau+2^e)$ for the integers $e\ge0$ with $2^e<H_\tau$.
\item The \emph{restart points} are the passages of all tops, and the starts $L_\tau(y_\tau)$ of the canonical excursions of
all tops except those that equal a descent-ladder time $D_\sigma(w)$ of a top~$\sigma$ at a level~$w$ of~$\sigma$.
\item The frame system of \emph{class~U} of~$x$ has the horizon~$T$, the walk~$x$, the frames $[0,T)$ and the suspend frames,
and these restart points; it is a frame system by \cref{lem:excursions}(c). The frame system of \emph{class~L} of~$x$ is that
of class~U of~$-x$. The upper hops of~$x$ belong to class~U, its lower hops to class~L.
\end{enumerate}
\end{definition}

The starts excluded in~(c) are those of the first excursions of descent structures that begin with their structure, at a
ladder time of the parent; their segments start one step later, at the start of the descent frame. Inside $E_\tau$, the own
times of the frame $\inner(\tau)$ are the times in no frame of~$\tau$ or of a descendant, among them the ladder times, and its
restarts are the passages (see the outline of \cref{lem:ladderbands}): the climb, the top and the descent down to the highest
passage form one segment, with home $y_\tau$ or $y_\tau+1$.

\paragraph{Why passages at powers of two.} Every suspend frame must end at band~$0$ (\cref{lem:W2b} below; the display needs
this), and its last time typically lies at the end of the descent of the frame's last excursion. Without a restart on the descent, the
running maximum reached at the top would stay until the excursion ends. The lowest passage, $D_\tau(y_\tau+1)$, restarts the
segment one level above the base; it is the passage that the proof of \cref{lem:W2b} uses. Every restart lowers the bands at the
ends of the rungs below it, and each lowered band costs its deficit below $B=\band(H)$. Restarting at the levels $y_\tau+2^e$
lets the band on the descent come down with the remaining height, as the band on the climb goes up with the height reached, and
it keeps the cost linear: the rung at $y_\tau+d$ with $2^{e-1}<d\le2^e<H_\tau$ ends at a band of at least $\band(2^e-d)$, and
the deficits $B-\band(2^e-d)$ over this dyadic block add up to $2^{e-1}(B-e+3)-2$ (\cref{lem:perexc}(a)), a series that is
geometric in $B-e$.
Restarts at every level of the descent would instead cost about~$B$ per rung, about $H\log_2H$ on a single excursion.

To state the bands at the ladder times, let $\lceil\log_2d\rceil$ be the least~$e$ with $d\le2^e$, and for $1\le d<H$ let
\[
  \mathrm{ub}(H,d)=\begin{cases}\band\bigl(2^{\lceil\log_2d\rceil}-d\bigr)&\text{if }2^{\lceil\log_2d\rceil}\le H-1,\\
  \band(H-1)&\text{otherwise}\end{cases}
\]
(\lean{uband}). The first case concerns the rungs at and below the \emph{top passage}, the highest passage of the excursion.

\begin{lemma}[Bands at ladder times]\label{lem:ladderbands}
\formal{LadderBands.lean}{home_mem_exc, beta_le_band_height, beta_ladL_bounds, beta_ladD_bounds, ladL_lt_segStart_ladD}
Let $x$ be a walk, take the frame system of class~U of~$x$, and let $\tau$ be a top with $E_\tau=(a,b)$, base~$y$, height~$H$
and $B=\band(H)$.
\begin{enumerate}[label=(\alph*)]
\item \emph{(Homes.)} For $a<t<b$, $y\le\home(t)\le x(\tau)$ and $\beta(t)\le B$.
\item \emph{(The climb.)} For $1\le d<H$, $\band(d-1)\le\beta(L_\tau(y+d))\le B$.
\item \emph{(The descent.)} For $1\le d<H$, $\mathrm{ub}(H,d)\le\beta(D_\tau(y+d))\le B$.
\item \emph{(Restarted rungs.)} If $1\le d<H$ and $2^{\lceil\log_2d\rceil}<H$, then the segment of the innermost frame at
$D_\tau(y+d)$ starts after $L_\tau(y+d)$.
\end{enumerate}
\end{lemma}

\begin{outline}
Let $F=\inner(\tau)$. Call a time \emph{free} if no frame of~$\tau$ or of a top whose excursion lies in~$E_\tau$ contains it
(\lean{Free}, \lean{frame_cases}); ladder times are free (\cref{lem:excursions}(d)). For $a<t<b$, $t$ is an own time of~$F$ if
and only if it is free (\lean{own_of_free}, \lean{free_of_own}), and the free restart points are the passages
(\lean{restart_inside}, \lean{passage_spec}). The excursion's segment starts at $\max(f_1,a)$, at the cell $y$ or $y+1$: $F$
starts at or before~$a$, or it is a descent frame starting at $a+1$ (\lean{inner_fst_cases}, \lean{excStart_spec},
\lean{x_excStart}). It is the segment of~$F$ at the climb ladder times, at~$\tau$, and at the descent ladder times above the top
passage (\lean{segStart_eq_excStart}, \lean{segStart_ladD_of_not}); below it, the segment at $D_\tau(y+d)$ starts at the passage
$D_\tau(y+2^{\lceil\log_2d\rceil})$ (\lean{segStart_ladD}). (a) Strictly inside the excursion every segment starts in $[a,t]$
(\lean{segStart_inner_mem}), so homes and own times lie in $[y,x(\tau)]$ (\lean{runMax_le_height}). (b) At $L_\tau(y+d)$ the
walk is at distance at least $d-1$ from the home (\cref{lem:dynamics}(e)). (c) Below the top passage the home is
$y+2^{\lceil\log_2d\rceil}$; above it the segment contains the own time~$\tau$, at distance at least $H-1$ from the home.
(d) The passage follows~$\tau$.
\end{outline}

\begin{lemma}[Lemma W2: suspend frames end at band~0]\label{lem:W2b}
\formal{LadderBands.lean}{beta_last_eq_zero, runMax_last_le_one, struct_exc, own_of_le_level}
Let $x$ be a walk, and take the frame system of class~U of~$x$. For every nonempty suspend frame $[f_1,f_2)$,
$\beta(f_2-1)=0$.
\end{lemma}

\begin{outline}
It suffices that the running maximum at $\ell=f_2-1$ is at most~$1$. Let $w$ be the level of $J=[f_1,f_2)$; on~$J$ the walk is at
least~$w$, and $x(f_2)\le w$ (\lean{frameAt_walk}), so $x(\ell)\in\{w,w+1\}$. Times of~$J$ at the level~$w$ are own times of~$J$
(\lean{own_of_le_level}), and the upper hops at the cell~$w$ in the structure are canonical excursions with frame~$J$
(\lean{struct_exc}). If $x(\ell)=w$ (\lean{runMax_last_le_one}), either the walk stays at~$w$ on~$J$, or the segment at~$\ell$
started at the lowest passage $D_E(w+1)$ of the last excursion~$E$ of the structure, or, if $H_E=1$, at the start of~$E$ or
one step later; as the descent structures of~$E$ are suspended, the own times of the segment are at the levels $w$ and $w+1$. If $x(\ell)=w+1$, the last
excursion~$E$ ends at~$f_2$, and $\ell$ is a time of~$E$ of height~$1$, the restart $D_E(w+1)$, or the last time of the descent
frame of~$E$ at level $w+1$, to which the first case applies.
\end{outline}

Lemma~W2 is the hypothesis \lean{FrameSystem.PopsAtZero} of view invariance (\cref{prop:W}), for the frame systems of all heads
(\lean{popsAtZero_headFrames}).

\subsection{Frame weight}\label{sec:ladders-weight}

The label counts segment starts as events, weighted by $2^{\text{scale}}$, where the scale of a segment is the band of its
maximum. Lemma~W1 bounds the total weight of each class of each walk by~$T$.

\begin{lemma}[Lemma W1: frame weight]\label{lem:W1}
\formal{FrameWeight.lean}{sup_add_one_le_card, FrameSystem.segOf, FrameSystem.segments, FrameSystem.segTimes, FrameSystem.segMax, FrameSystem.sum_segWeight_le}
\formal{FrameWeightU.lean}{own_step, exists_outer_subframe, own_frame_end, segMax_add_one_le, sum_segWeight_upperFrames_le, card_levels, card_framesOf_le, card_suspendFrames_le}
Let $\mathcal S$ be a frame system. The \emph{segment} of a time $t<T$ is $(\inner(t),s_{\inner(t)}(t))$, and the segments
of~$\mathcal S$ are the segments of these times. The \emph{own times} of a segment $(F,s)$ are the times $t<T$ whose segment it
is, and its \emph{maximum} $m(F,s)$ is the largest $\abs{x(t)-x(s)}$ over its own times ($0$ if there is none).
\begin{enumerate}[label=(\alph*)]
\item If every segment $(F,s)$ has $m(F,s)$ smaller than its number of own times, then the weights $2^{\band(m(F,s))}$ of all
segments add up to at most~$T$.
\item The frame system of class~U of a walk satisfies the hypothesis of~(a), so its segment weights add up to at most~$T$.
\item A walk has at most $2T$ suspend frames.
\end{enumerate}
\end{lemma}

\begin{outline}
(a) The segments partition $[0,T)$, and $2^{\band(m)}\le m+1$. (b) The \emph{staircase} (\lean{sup_add_one_le_card}): a function
that vanishes at the least element of a finite set~$A$, and at every other element is at most one above its value at some
earlier element, has maximum at most $\card A-1$. Consecutive own times of a frame are at most one level apart
(\lean{own_step}): they are consecutive, or the times between them form a frame of a top whose excursion lies directly in the
frame (\lean{exists_outer_subframe}), and they are the rungs $L(w-1),L(w)$ or $D(w),D(w-1)$ around it (\lean{x_frame_ends},
\lean{own_frame_end}). The first own time of a segment is at its home: it is~$s$, or, if $s$ is the start of~$F$ but not an own
time, the end of the outermost hanging frame that starts there, where the walk is again at its level (\lean{x_hangFrame_fst}).
(c) A top of height~$H$ has $H-1$ levels (\lean{card_levels}), hence at most $2(H-1)$ frames (\lean{card_framesOf_le}), and the
heights add up to at most~$T$ (\cref{lem:ladder}(d)).
\end{outline}

\subsection{Kept and lost hops}\label{sec:ladders-kept}

The following definitions apply to every frame system~$\mathcal S$ with walk~$x$ and horizon~$T$, and every hop $e=(p,u)$
of~$x$, with owner~$\tau$ and cell $z=x(p)$. In Theorem~A the frame system is that of the hop's class.

\begin{definition}[Kept and lost hops]\label{def:kept}
\formal{HopData.lean}{bandP, bandU, ownerScale, homeBand, segScale, pushMax, FrameMicro, Lost, cstar}
\formal{CostModel.lean}{Aligned, Ground, hopCost}
The \emph{data} of the hop are the bands $\beta_p=\beta(p)$ and $\beta_u=\beta(u)$; the \emph{owner scale}
$\sigma_e=\band(\prom(\tau))$; the \emph{home band} $\lambda_e=\band(S_e)$, where $S_e=\max_{p\le c\le u}\abs{z-\home(c)}$; the
largest scale $\pi_e$ of a segment that starts in $(p,u]$, that is, the maximum of $\band(m(F,s))$ over the frames~$F$ and the
times $s\in(p,u]$ that are the start or a restart of~$F$ ($0$ if there are none); and whether the hop is \emph{frame-micro}, that
is, whether the segment of~$u$ started at or before~$p$: $s_{\inner(u)}(u)\le p$. Let $h_0\ge1$ and $r<h_0$, with $G$
and~$\theta$ as in \cref{def:grid}. The hop is
\begin{itemize}[leftmargin=1.5em]
\item \emph{aligned} at~$r$ if $G(\beta_p)=G(\beta_u)=:d$, $\lambda_e<\theta(d+1)+2$ (it is \emph{home-near}:
$S_e<2^{\theta(d+1)+2}$), and it is not frame-micro or $\theta(d)\le\sigma_e$ or $\theta(d)\le\pi_e$ (it is \emph{eventful});
\item \emph{ground} at~$r$ if $G(\sigma_e)=0$;
\item \emph{kept} at~$r$ if it is aligned or ground, and \emph{lost} at~$r$ otherwise.
\end{itemize}
Its \emph{cost} is
\[
  c_e=\bigl(\max(\beta_p,\beta_u)-\min(\beta_p,\beta_u)\bigr)+\bigl(\lambda_e-2-\max(\beta_p,\beta_u)\bigr)^+
  +[\text{frame-micro}]\cdot\bigl(\beta_u-\max(\sigma_e,\pi_e)\bigr)^+,
\]
where $a^+=\max(a,0)$ and $[\cdot]\in\{0,1\}$, and its \emph{capped cost} is $c^*_e=\min(c_e,\sigma_e)$.
\end{definition}

By Lemma~V (\cref{lem:V}), a kept hop is cut only if the display is truncated, changes above the hop's deciding group~$d$,
or a digit wraps. In one group, both ends read the same digit; home-nearness places the cell within $2^{\theta(d+1)+2}$ of every
home active during the hop, which bounds the number of kept hops that span a time (\cref{sec:accounting}); eventfulness provides
an event of scale at least $\theta(d)$ inside the hop (\cref{lem:validity}). A ground hop has extent less than $2^{\theta(1)}$, by
\cref{prop:prom} (\lean{sub_le_prom}). Lost hops are cut, and Theorem~A bounds their number.

\begin{lemma}[The cost model]\label{lem:cost}
\formal{CostModel.lean}{card_lost_le, hopCost_le_of_bands}
\formal{HopData.lean}{sum_card_lost_le_sum_cstar, sum_cstar_le_of_ladders}
Let $h_0\in\N$.
\begin{enumerate}[label=(\alph*)]
\item For all values of the data $(\beta_p,\beta_u,\sigma,\lambda,\pi)$ and of frame-micro, the hop is lost for at most
$\min(c,\sigma)$ of the offsets $r<h_0$, where $c$ is its cost.
\item If $\beta_p,\beta_u,\lambda\le B$, and $\beta_u\le\max(\sigma,\pi)$ in case the hop is frame-micro, then
$c\le B-\min(\beta_p,\beta_u)$.
\item For every frame system and every finite set~$\mathcal H$ of hops,
$\sum_{r<h_0}\#\{e\in\mathcal H:e\text{ lost at }r\}\le\sum_{e\in\mathcal H}c^*_e$.
\item Let $f$ be a function on hops and $C\in\N$. If $\sum_{e\in\text{ladder}(\tau)}f(e)\le C\,H_\tau$ for every top~$\tau$ of a
map $x\colon\N\to\Z$, then $\sum_ef(e)\le C\cdot\#\{\text{upper hops of }x\}$, the sum running over the upper hops of~$x$.
\end{enumerate}
\end{lemma}

\begin{outline}
(a) A union bound with \cref{lem:grid}. A hop that is neither aligned nor ground has differing groups at its ends, for at most
$\max-\min$ offsets; or equal groups and $\theta(G(\max(\beta_p,\beta_u))+1)\le\lambda-2$, for at most $(\lambda-2-\max)^+$
offsets; or it is frame-micro with $\max(\sigma,\pi)<\theta(G(\beta_u))$, for at most $(\beta_u-\max(\sigma,\pi))^+$ offsets.
And a hop that is not ground has $G(\sigma)\ge1$, for at most~$\sigma$ offsets. (b) The last term vanishes, and
$(\lambda-2-\max)^+\le B-\max$; so $c\le B-\min$. (c) Exchange the sums and apply~(a). (d) Group the upper hops by their tops
and use $\sum_\tau H_\tau=\#\{\text{upper hops}\}$ (\lean{sum_height_eq}).
\end{outline}

So, for a uniform offset, a hop is lost with probability at most $c^*_e/h_0$. The cap~$\sigma_e$ matters for hops whose owner is
not prominent: their ends may have unrelated bands, but they can be lost only if a grid floor lies at or below the owner's scale.

\paragraph{Events.} The \emph{events} of a frame system are the \emph{turning points}, the times $s<T$ with $\prom(s)\ge1$, of
scale $\band(\prom(s))$, and the \emph{segment starts}, the pairs $(F,s)$ of a frame~$F$ and a time $s<T$ that is the start or a
restart of~$F$, of scale $\band(m(F,s))$. $N_\theta(t)$ is the number of events at times at most~$t$ of scale at least~$\theta$
(\lean{eventCount}); it is nondecreasing in~$t$ (\lean{eventCount_mono}). The \emph{deciding group} of a hop at the offset~$r$
is $G(\beta_u)$ if the hop is aligned at~$r$, and $0$ otherwise (\lean{decGroup}).

\begin{lemma}[Validity]\label{lem:validity}
\formal{Events.lean}{turnCount, segStarts, segEventCount, eventCount, eventCount_mono, decGroup, owner_turn, eventCount_lt_decGroup}
Let $\mathcal S$ be a frame system with walk~$x$ and horizon~$T$, let $(p,u)$ be a hop of~$x$ that is an upper hop, and let
$h_0,r\in\N$. If $d$ is the deciding group of the hop at~$r$, then $N_{\theta(d)}(p)<N_{\theta(d)}(u)$.
\end{lemma}

\begin{outline}
The owner~$\tau\in(p,u)$ has $\prom(\tau)\ge x(\tau)-x(p)\ge1$ by \cref{prop:prom} (\lean{owner_turn}), so it is a turning point
of scale~$\sigma_e$. If the hop is not aligned, $d=0$ and $\theta(0)=0\le\sigma_e$. If it is aligned, eventfulness provides the
event: if it is not frame-micro, the segment of~$u$ starts in $(p,u]$, a segment start of scale at least
$\beta(u)\ge\theta(G(\beta_u))$ (\lean{beta_le_segScale}, \cref{lem:grid}); if $\theta(d)\le\sigma_e$, the owner serves; if
$\theta(d)\le\pi_e$, the owner serves when $\theta(d)=0$, and otherwise a segment start in $(p,u]$ of scale $\pi_e\ge\theta(d)$.
\end{outline}

Lemma~V needs this only for kept hops; the formal statement holds for every upper hop, since a hop that is not aligned has
deciding group~$0$. For class~L it applies to the frame system of~$-x$.

\subsection{The per-excursion bound}\label{sec:ladders-perexc}

We add up the capped costs along one ladder. The arithmetic is separate from the walk: the rung at the level $d\ge1$ above the
base costs at most the deficits $B-\band(d-1)$ and $B-\mathrm{ub}(H,d)$ of its ends, and the hop at the base at most~$B$.

\begin{lemma}[The ladder sum]\label{lem:perexc}
\formal{PerExcursion.lean}{uband, sum_range_two_pow_sub_band, sum_sub_band_two_pow_clog, ladder_sum_le_six, sum_range_le_six_of_ladder}
Subtraction of natural numbers is truncated at~$0$.
\begin{enumerate}[label=(\alph*)]
\item \emph{(A dyadic block.)} For $C,e\in\N$ with $e\le C+1$, $\sum_{w<2^e}(C-\band(w))+2=2^e(C+2-e)$.
\item \emph{(The descent up to a passage.)} For $C,e\in\N$ with $e\le C+1$,
$\sum_{d=1}^{2^e}\bigl(C-\band(2^{\lceil\log_2d\rceil}-d)\bigr)+2e+4=2^e(C+4-e)$.
\item \emph{(The ladder sum.)} For every $H\in\N$, with $B=\band(H)$,
\[
  B+\sum_{d=1}^{H-1}\Bigl(\bigl(B-\band(d-1)\bigr)+\bigl(B-\mathrm{ub}(H,d)\bigr)\Bigr)\le6H.
\]
\item Consequently, if $c\colon\N\to\N$ satisfies $c(0)\le B$ and $c(d)\le(B-\band(d-1))+(B-\mathrm{ub}(H,d))$ for $1\le
d<H$, then $\sum_{d<H}c(d)\le6H$.
\end{enumerate}
\end{lemma}

\begin{outline}
(a) Induction on~$e$: the levels $w\in[2^e,2^{e+1})$ have band~$e$ (\lean{band_eq_of_mem_Ico}). (b) Induction on~$e$: the levels
$d\in(2^e,2^{e+1}]$ have $\lceil\log_2d\rceil=e+1$ (\lean{clog_eq_of_lt_of_le}), and $w=2^{e+1}-d$ maps them onto the block of~(a).
(c) For $H\le1$ by inspection. Otherwise let $e=\band(H-1)$, so that $2^e\le H-1<2^{e+1}$, $B\in\{e,e+1\}$, and the top passage
is at the level $2^e$. The climb side is at most $\sum_{j=0}^{H}(B-\band(j))=2^{B+1}-2$, by Lemma~E (\lean{sum_band_sub_eq},
\cref{lem:bandsums}). The descent side contributes $2^e(B+4-e)-2e-4$ for $d\le2^e$, by~(b) with $C=B$, and at most $1$ for each
$d\in(2^e,H)$, where $\mathrm{ub}(H,d)=\band(H-1)\ge B-1$. If $B=e$, the total is at most $5\cdot2^e+H-e-7\le6H$; if $B=e+1$, then
$H=2^{e+1}$ and the total is at most $5H-e-6$. (d) Apply~(c).
\end{outline}

\begin{lemma}[The per-excursion bound]\label{lem:perexc-bound}
\formal{TheoremA.lean}{ownerScale_ladderHop, cstar_exc_le, cstar_ladder_le, sum_ladder_cstar_le}
Let $x$ be a walk, take the frame system of class~U of~$x$, and let $\tau$ be a top of height~$H$ and base~$y$, with
$B=\band(H)$. Every hop of the ladder of~$\tau$ has owner scale $\band(\prom(\tau))$, and:
\begin{enumerate}[label=(\alph*)]
\item \emph{(The base hop.)} $c^*_{E_\tau}\le B$.
\item \emph{(The rungs.)} For $1\le d<H$, the hop of the ladder at the cell $y+d$ has $c^*\le(B-\band(d-1))+(B-\mathrm{ub}(H,d))$.
\item \emph{(The ladder.)} $\sum_{e\in\text{ladder}(\tau)}c^*_e\le6H$.
\end{enumerate}
\end{lemma}

\begin{outline}
(b) The owner is~$\tau$, so $\sigma\ge B$ by \cref{prop:prom} (\lean{height_le_prom}). By \cref{lem:ladderbands}, the bands at the
ends lie between $\band(d-1)$, respectively $\mathrm{ub}(H,d)$, and~$B$, and the home band is at most~$B$. The hop is eventful:
below the top passage it is not frame-micro (\cref{lem:ladderbands}(d)), and above it $\beta_u\le B\le\sigma$. So
$c\le B-\min(\beta_p,\beta_u)$ by \cref{lem:cost}(b). (a) If $E_\tau$ is the first excursion of a descent structure of a top~$P$,
then $P$ is earlier and at least as high as~$\tau$; if it is the last excursion of a hanging structure of~$P$, then $P$ is later
and strictly higher (\lean{inner_fst_cases}, \lean{inner_snd_cases}). In both cases the walk stays at or above~$y$ in between, so
$\prom(\tau)\le H$ by \cref{prop:prom} (\lean{prom_le_of_left}, \lean{prom_le_of_right}) and $c^*\le\sigma\le B$: this is where
the cap is needed, as one end of the hop lies in the parent's frame. Otherwise the segments at the times of $[a,b]$ start in the
excursion (\lean{segStart_inner_mem_Icc}), the bands and the home band are at most~$B$ (\lean{runMax_le_height}), and the hop is
frame-micro only if $H=1$, when $\beta_u=0$: for $H\ge2$ the passage $D_\tau(y+1)$ restarts the segment of~$b$ after~$a$. (c) The
hops at the cells $y+d$, $d<H$, are the ladder (\cref{lem:ladder}(a)); apply \cref{lem:perexc}(d).
\end{outline}

\subsection{Theorem A}\label{sec:ladders-thmA}

\begin{theorem}[Theorem A: lost hops of one walk]\label{thm:A}
\formal{TheoremA.lean}{thmA_upper, thmA_lower}
Let $x$ be a walk, $T$ a horizon and $h_0\in\N$. Then
\begin{align*}
  \sum_{r=0}^{h_0-1}\#\{\text{upper hops of $x$ lost at $r$ in class U}\}&\le6\cdot\#\{\text{upper hops of $x$}\},\\
  \sum_{r=0}^{h_0-1}\#\{\text{lower hops of $x$ lost at $r$ in class L}\}&\le6\cdot\#\{\text{lower hops of $x$}\}.
\end{align*}
That is, for $h_0\ge1$ and a uniformly random offset $r\in\{0,\dots,h_0-1\}$, the expected number of lost hops of each class is
at most $6/h_0$ times the number of hops of that class.
\end{theorem}

\begin{outline}
For class~U, \cref{lem:cost}(c) bounds the left side by $\sum_ec^*_e$ over the upper hops, and \cref{lem:cost}(d) with
\cref{lem:perexc-bound}(c) bounds this by $6$ times their number (\lean{thmA_upper}). For class~L, apply this to the walk~$-x$,
whose upper hops are the lower hops of~$x$ (\cref{lem:ladder}(e), \lean{thmA_lower}).
\end{outline}

Since a walk has at most $T$ hops (\cref{sec:prelim-walks}), the expected number of lost hops of its two classes together is at
most $6T/h_0$; the proof of the main theorem chooses $h_0$ so large that this is a small fraction of~$T$ (\cref{sec:proof}).
Lemmas~W1 and~W2 are the other two inputs from this section; the kept hops are the subject of the following sections.

\begin{remark}[No pause-freeness is needed]\label{rem:A-pauses}
Every result of this section assumes only that $x$ is a walk; pauses are allowed. A walk may pause at its top
(\cref{ex:ladders}), and then the own times of an excursion's frame are not all ladder times; the proofs use only own times,
ladder times and passages. Since the composition of the following sections does not use pause-freeness either, the main theorem
is proved for walks with pauses directly, at the depth of $k$ walks (\lean{tcpAt_direct}), with no detour through pause-free
walks.
\end{remark}
